\chapter{面向对象}

\section{类与对象}
\subsection{类}
定义一个类，使用关键字\texttt{class}，类名的命名规则与变量名相同，首字母大写。

\begin{lstlisting}[style=demostyle,numbers=none]
class Person[(father)]: # 类首字母大写 [] 部分可选，表示继承的父类
    def __init__(self, name, age):
        self.name = name
        self.age = age

    def say(self):
        print('hello world')

p = Person('Tom', 20) # 实例化
print(p.name)
print(p.age)
p.say()
\end{lstlisting}

\subsection{继承}
Python支持多继承，使用逗号分隔。
\begin{lstlisting}[style=demostyle,numbers=none]
class Teacher:
    def teaching(self):
        print('I am teaching')

class Dad:
    def cooking(self):
        print('I am cooking')

class Son(Teacher, Dad):
    def playing(self):
        print('I am playing')

s = Son()
s.teaching()
s.cooking()
s.playing()
\end{lstlisting}

\subsection{接口}
Python中没有接口的概念，但是可以使用抽象类来实现接口的功能。
\begin{lstlisting}[style=demostyle,numbers=none]
class Animal:
    def run(self):
        pass # 占位符

class Dog(Animal):
    def run(self):
        print('dog is running')

class Cat(Animal):
    def run(self):
        print('cat is running')

d = Dog()
d.run()

c = Cat()
c.run()
\end{lstlisting}


\subsection{封装}
\subsubsection{私有属性}
Python中没有私有属性的概念，但是可以通过在属性名前加两个下划线来实现私有属性。
\begin{lstlisting}[style=demostyle,numbers=none]
class Person:
    def __init__(self, name, age):
        self.__name = name
        self.__age = age

    def say(self):
        print('hello world')

p = Person('Tom', 20)
print(p.__name) # 不能访问

p.say()
\end{lstlisting}

\subsubsection{私有方法}
Python中没有私有方法的概念，但是可以通过在方法名前加一个下划线来实现私有方法。
\begin{lstlisting}[style=demostyle,numbers=none]
class Person:
    def __init__(self, name, age):
        self.__name = name
        self.__age = age

    def __say(self):
        print('hello world')

    def say(self):
        self.__say()     

p = Person('Tom', 20)
p.say()
\end{lstlisting}

\subsection{多态}
多态是指不同的对象对同一消息做出不同的响应。
\begin{lstlisting}[style=demostyle,numbers=none]
class Animal:
    def run(self):
        pass

class Dog(Animal):
    def run(self):
        print('dog is running')

class Cat(Animal):
    def run(self):
        print('cat is running')

def run(animal):
    animal.run()

d = Dog()
run(d)

c = Cat()
run(c)
\end{lstlisting}




